\subsection{Definition and examples}

Recall (Ch. V, §6, Nos. 1 and 4; Ch. VI, §2, No.~10) that, if X and Y are locally compact spaces, \(\mu\) a measure on X, and \(\varphi\) a mapping of X into Y, \(\varphi\) is said to be \(\mu\) -proper if: \(a)\) \(\varphi\) is \(\mu\) -measurable; \(b)\) for every compact subset K of Y, \(\overline{\varphi}^1 (\mathbf{K})\) is essentially \(\mu\) -integrable. Then the image measure \(\nu = \varphi (\mu)\) on Y exists and has the following property: for a function \(\mathbf{f}\) on Y, with values in a Banach space or in \(\overline{\mathbf{R}}\) , to be essentially integrable for \(\nu\) , it is necessary and sufficient that \(\mathbf{f} \circ \varphi\) be so for \(\mu\) , in which case,

\[
\int_ {\mathrm{Y}} \mathbf {f} (y) d \nu (y) = \int_ {\mathrm{X}} \mathbf {f} (\varphi (x)) d \mu (x).
\]

DEFINITION 1. --- Let \(X_1, \ldots, X_n\) be locally compact spaces, \(\mu_i\) a measure on \(X_i\) ( \(1 \leqslant i \leqslant n\) ); let \(X\) be the product of the \(X_i\) , \(\mu\) that of the \(\mu_i\) . Let \(\varphi\) be a mapping of \(X\) into a locally compact space \(Y\) . One says that the sequence \((\mu_i)\) is \(\varphi\) -convolvable, or that \(\mu_1, \ldots, \mu_n\) are \(\varphi\) -convolvable, if \(\varphi\) is \(\mu\) -proper; in this case, the image \(\nu = \varphi(\mu)\) of \(\mu\) under \(\varphi\) is called the convolution product of the \(\mu_i\) for \(\varphi\) , and is denoted \(*^{\varphi(\mu_i)}_{1 \leqslant i \leqslant n}\) , or \(\underset{i=1}{\overset{n}{*}} \mu_i\) , or \(\mu_1 * \mu_2 * \cdots * \mu_n\) .

The last two notations are of course used only when there can be no doubt as to \(\varphi\) .

Let f be a function on Y, with values in a Banach space or in \(\overline{\mathbf{R}}\) . In order that f be essentially integrable for \(\mu_{1}*\cdots*\mu_{n}\) , it is necessary and sufficient that the function

\[
\left(x _ {1}, \dots , x _ {n}\right) \mapsto \mathbf {f} \left(\varphi \left(x _ {1}, \dots , x _ {n}\right)\right)
\]

be essentially integrable for \(\mu_{1}\otimes\mu_{2}\otimes\cdots\otimes\mu_{n}\) , in which case

\[
\int \mathbf {f} d (\mu_ {1} * \dots * \mu_ {n}) = \int \mathbf {f} \big (\varphi (x _ {1}, \ldots , x _ {n}) \big) d \mu_ {1} (x _ {1}) \ldots d \mu_ {n} (x _ {n}),\tag{1}
\]

a formula that may be regarded as defining \(\mu_{1}*\cdots *\mu_{n}\) when one takes \(\mathbf{f}\in\mathcal{K}(\mathrm{Y})\) .

The definitions imply at once that the \(\mu_{i}\) are convolvable if and only if the \(|\mu_{i}|\) are. When this is the case,

\[
\left| \varphi \left(\mu_ {1} \otimes \dots \otimes \mu_ {n}\right) \right| \leqslant \varphi \left(\left| \mu_ {1} \otimes \dots \otimes \mu_ {n} \right|\right) = \varphi \left(\left| \mu_ {1} \right| \otimes \dots \otimes \left| \mu_ {n} \right|\right)
\]

(Ch. VI, §2, No.~10), that is,

\[
\left| \begin{array}{c} * \\ i \end{array} \mu_ {i} \right| \leqslant * _ {i} | \mu_ {i} |.\tag{2}
\]

If the \(\mu_{i}\) are convolvable and positive, and if \(\nu_{i}\) is a measure on \(X_{i}\) such that \(0 \leqslant \nu_{i} \leqslant \mu_{i}\) , then the \(\nu_{i}\) are convolvable and

\[
\underset {i} {\ast} \nu_ {i} \leqslant \underset {i} {\ast} \mu_ {i}.
\]

Suppose \(\mu_1, \mu_2, \ldots, \mu_n\) are convolvable, and that \(\mu_1', \mu_2, \ldots, \mu_n\) are convolvable ( \(\mu_1'\) being a measure on \(X_1\) ). By Ch. V, §6, No.~3, Cor. 1 of Prop. 6, \(\mu_1 + \mu_1', \mu_2, \ldots, \mu_n\) are convolvable and

\[
\left(\mu_ {1} + \mu_ {1} ^ {\prime}\right) * \mu_ {2} * \dots * \mu_ {n} = \mu_ {1} * \mu_ {2} * \dots * \mu_ {n} + \mu_ {1} ^ {\prime} * \mu_ {2} * \dots * \mu_ {n}.
\]

Examples. --- 1) For any \(\varphi\) , the measures \(\varepsilon_{x_i}\) , where \(x_i \in X_i\) for \(1 \leqslant i \leqslant n\) , are always convolvable and have convolution product \(\varepsilon_y\) , with \(y = \varphi(x_1, x_2, \ldots, x_n)\) . Consequently, if each of the \(\mu_i\) has finite support, then the \(\mu_i\) are convolvable and \(\mu_1 * \cdots * \mu_n\) has finite support. In particular, let M be a monoid \(^1\) equipped with a locally compact topology; if one takes \(\varphi\) to be the law of composition in M then the measures on M with finite support form, for convolution, an algebra that is none other than the algebra of the monoid M (over R or over C, according as one considers real or complex measures) (A, III, §2, No.~6).

\begin{enumerate}
\def\labelenumi{\arabic{enumi})}
\setcounter{enumi}{1}
\tightlist
\item
  Let M be a monoid equipped with the discrete topology; assume that for each \(m \in M\) , there are only finitely many pairs \((m', m'') \in M \times M\) such that \(m'm'' = m\) ; this amounts to saying that the law of composition in M is a proper mapping of \(M \times M\) into M; the measures on M then form an algebra for convolution, an algebra that is none other than the total algebra of the monoid M (A, III, §2, No.~10); we note the following two special cases:
\end{enumerate}

\begin{enumerate}
\def\labelenumi{\alph{enumi})}
\tightlist
\item
  M = N, the law of composition being addition. To every measure \(\mu\) on N, let us associate the formal series
\end{enumerate}

\[
\mathrm{S} (\mu) = \sum_ {n = 0} ^ {\infty} \mu (\{n \}) t ^ {n}
\]

in an indeterminate t. Then \(S(\mu * \mu') = S(\mu)S(\mu')\) . An analogous remark holds for formal series in any number of indeterminates.

\(^{*}b)\) M = N \(^{*}\) , the law of composition being multiplication. To every measure \(\mu\) on N \(^{*}\) , let us associate the formal Dirichlet series

\[
\mathrm{D} (\mu) = \sum_ {n = 1} ^ {\infty} \mu (\{n \}) n ^ {- s}.
\]

Then \(\mathrm{D}(\mu * \mu') = \mathrm{D}(\mu) \mathrm{D}(\mu')\) .

\begin{enumerate}
\def\labelenumi{\arabic{enumi})}
\setcounter{enumi}{2}
\tightlist
\item
  Let X,Y,Z be locally compact spaces, \(\varphi\) a continuous mapping of \(X \times Y\) into Z. If \(x \in X\) and \(\mu\) is a measure on Y, to say that \(\varepsilon_{x}\) and \(\mu\) are \(\varphi\) -convolvable comes to saying that the mapping \(\varphi(x,\cdot)\) of Y into Z is \(\mu\) -proper. One then has \(\varepsilon_{x} * \mu = \varphi(x,\cdot)(\mu)\) .
\end{enumerate}

